% Midterm Exam -- Spring 2026
% ECO 4971/6971 -- Statistical Learning (ISLR)
% Material: Classes 5, 6, 7, 8, 8B

\documentclass[11pt]{article}
\usepackage[margin=1in]{geometry}
\usepackage{amsmath,amssymb}
\usepackage{enumitem}
\usepackage{booktabs}
\usepackage{graphicx}
\usepackage{float}

\title{Midterm Exam -- Spring 2026\\[0.5em]
  \large ECO 4971/6971 -- Machine Learning and Econometrics}

\begin{document}
\maketitle
\noindent \textbf{Instructions:} Short-answer, pen-and-paper. Show brief reasoning where appropriate.

\bigskip
\hrule
\bigskip



\begin{enumerate}[leftmargin=*]

\item In OLS, the Residual Sum of Squares (RSS) is defined as
$$
  \text{RSS} = \sum_{i=1}^n (y_i - \hat{y}_i)^2.
$$

\begin{enumerate}
  \item How is RSS related to MSE? 
  \vspace{2cm}  
  
  \item Explain why squaring the errors (rather than taking absolute values) to penalize some kinds of errors more heavily than others. What kind of error does RSS penalize the most?
  \vspace{4cm}  
  \item If two models are fit to the same data, and Model 1 has a much smaller RSS than Model 2, what (if anything) can you conclude about their predictive performance on new data?
  \vspace{3cm}  
\end{enumerate}
 
\clearpage

\item This question covers bias and variance as taught in this class (not your econometrics class). Consider two KNN regression models on the same dataset: Model A uses $K=3$, and Model B uses $K=25$. Answer the following questions:  
\begin{enumerate}
  \item In one sentence each, explain what bias and variance are in the context of this class. 
  \vspace{3cm}  
  \item Which of the two KNN models will have a higher bias and which will have a higher variance? You don't need to explain. 
  \vspace{1cm}  
  \item Which model will likely have a higher training MSE and which will likely have a higher test MSE? Explain briefly. 
  \vspace{3cm}  
\end{enumerate}



\item In the linear probability model (LPM), we regress a binary outcome $Y$ on predictors using OLS. Why are Economists more likely to use LPM, but we tend to use other models in a statistical learning context? Be sure to include an advantage (for econometrics) and a disadvantage (for statistical learning) of a LPM. 

\clearpage


\item  You are shown a table of predicted probabilities from a logistic regression model for 10 observations, along with the true labels (0 or 1). At a threshold of 0.5, the confusion matrix is:
\begin{center}
\begin{tabular}{l|cc}
 & Predicted 0 & Predicted 1 \\\hline
Actual 0 & 6 & 1 \\
Actual 1 & 1 & 2 \\
\end{tabular}
\end{center}
If the threshold is lowered to 0.3, the confusion matrix becomes:
\begin{center}
\begin{tabular}{l|cc}
 & Predicted 0 & Predicted 1 \\\hline
Actual 0 & 5 & 2 \\
Actual 1 & 0 & 3 \\
\end{tabular}
\end{center}
\begin{enumerate}
  \item How did precision and recall for the positive class change when the threshold was lowered from 0.5 to 0.3?
  \vspace{3cm}  
  \item Explain briefly why this change is typical when we lower the threshold.
  \vspace{2cm}  
\end{enumerate} 


\item Below are four research questions. For each one, indicate whether it is more appropriate to use a \emph{machine-learning style} model taught in this calss or an \emph{econometrics style} model. Briefly justify your choice in one sentence per part.


\begin{enumerate}[label=\alph*)]
  \item Using past browsing behavior, can we predict whether a user will click on a particular online ad? 
  \vspace{1cm}    
  \item What is the effect of raising the minimum wage on employment among low-wage workers? 
  \vspace{1cm}  
  \item If a firm doubles its TV advertising budget, by how much will its sales change, holding all else equal?
  \vspace{1cm}  
  \item Given a customer’s credit history and demographics, what is the probability that they will default on their loan next month?
\end{enumerate}




\clearpage
\item This question is designed to test your ability to reason through an unfamiliar problem related to the ROC curve. 
\begin{enumerate}
  \item First, draw what you would expect a typical ROC curve for a generic classifier to look like, assuming the classifier is reasonably good. Be sure to label the axes. 
  \item Now, imagine this classifier is trained on two different datasets. In dataset \#1, the response is very rare (e.g., 1\% positives). In dataset \#2, there is a greater balance between positives and negatives (eg. 50\% positives). Both classifiers perform similarly: A model that classifies 90\% of the true positives correctly makes 5 false positives for every 10 true positives. Should the ROC curves look the same? If not, draw a second ROC curve for the second dataset and label it. 
\end{enumerate}

\end{enumerate}

\clearpage \newpage     

\begin{center}\textbf{Scratch Paper}\end{center}

\noindent
\vspace{16cm}

\end{document}

